\centerline{\bf \Huge Матметоды}
\title{Матметоды}

\begin{flushright}
    \scriptsize У людей с разными взглядами разные методы.\\
    Янь Фэй. 
\end{flushright}

\begin{flushleft}
  \textbf{\Large Введение переменной}
\end{flushleft}

\problem{1} Очир велел Расулу разделить число на 2, а к результату
прибавить 3. Расул же по ошибке умножил число на 2, а от полученного
произведения отнял 3. Но ответ все равно получился правильный. Какой?\\

\textbf{Ответ: 5.}\par

\textbf{Решение:} Мы бы легко решили задачу, если бы знали исходное число. Обозначим его символом x. Расул должен был вычислить $\dfrac{1}{2}\cdot x + 3$, а вычислил $2\cdot x - 3$. Раз он получил правильный ответ, то  $\dfrac{1}{2}\cdot x + 3 = 2\cdot x - 3$. Другими словами, $(3+3)=(2-\dfrac{1}{2})x$, т.е. $x=4$, а ответ 5.\\


\problem{2}Фигура на рисунке составлена из квадратов. Найдите сторону
левого нижнего квадрата, если известно, что сторона самого маленького
квадрата равна 1.

\begin{figure}[h] 
    \includegraphics[scale=1.5]{ЭЛМШ_2/picture.png}
\end{figure}

\textbf{Ответ: 4.}\par

\textbf{Решение:} Обозначим сторону среднего нижнего квадрата за x. Тогда другие квадраты имеют
размеры как на рисунке. Так как ширина картинки вверху и внизу одна и та же, $(x+3) + (x+2) = y+x +(x+1).$ Значит, $y = 5 - 1 = 4$\newpage


\begin{flushleft}
  \textbf{\Large Метод от противного}
\end{flushleft}

\problem{1}У Эрдема есть банки только с красной и синей краской. Банки с красной краской всегда большие. Есть маленькая банка с краской. Докажите, что краска в ней синяя.\\

\textbf{Решение:} Предположим, что краска в ней не синяя. Значит, краска внутри красная, но банки с красной краской всегда большие. Следовательно, она не может быть внутри маленькой банки.\\

\problem{2} Баир совершенно не любит задачи на скорость и не помнит ни одной формулы. Когда его спросили, какое расстояние проедет поезд, он попробовал и перемножить данные скорость и время, и сложить их, и даже поделить скорость на время. «У меня всегда получается одно и то же число! Наверное, это и есть правильный ответ!» — воскликнул Баир. Докажите, что выполнять арифметические действия Баир тоже не умеет.\\

\textbf{Решение:} Предположим противное, то есть предположим, что могут найтись такие скорость и время, что их сумма, произведение и результат деления скорости на время, равны.
Для начала взглянем на условие, что равны произведение скорости на время и результат деления скорости на время:
Если бы время было больше 1, то при умножении скорости на время, получилось бы число, большее данной скорости, а при делении скорости на время – меньшее;
Если бы время было меньше 1, то при умножении скорости на время, получилось бы число, меньшее данной скорости, а при делении скорости на время – большее.
Следовательно, если у Баира и могло так получится, что равны произведение скорости на время и результат деления скорости на время, то только в случае, когда время равно 1.
Но тогда точно будут не равны сумма скорости и времени: если умножить скорость на 1, то она не изменится, а если прибавить 1 – увеличится. Противоречие. Значит, наше предложение неверно и Баир выполнять арифметические операции не умеет.\\

\begin{flushleft}
  \textbf{\Large Принцип Дирихле}
\end{flushleft}

\problem{1}В поход пошли 25 туристов. Самому старшему из них 40 лет, а самому младшему – 17 лет. Верно ли, что среди туристов есть ровесники?
\newpage

\textbf{Ответ: Да, верно.}\par
\textbf{Решение:} Так как старшему 40 лет, а младшему 17, тогда давайте выпишем всевозможные варианты возрастов туристов: 17, 18, 19, 20, 21, 22, 23, 24, 25, 26, 27, 28, 29, 30, 31, 32, 33, 34, 35, 36, 37, 38, 39, 40.
Всего 24 различных варианта, конечно, мы могли посчитать это и по другому: 40-17+1= 24. Так как всего туристов 25, а различных вариантов возраста 24. Тогда пусть различные варианты - это клетки, а туристы - это кролики. По принципу Дирихле, так как количество кроликов больше количества клеток, следовательно, хотя бы в одной клетке будет как минимум 2 кролика, а это означает, что как минимум у двух туристов одинаковый возраст, то есть они ровесники\\

\problem{2}Алина загадала 8 различных натуральных чисел, не больших 15. Докажите, что среди их положительных попарных разностей Алина всегда найдет три одинаковых.\\

\textbf{Решение}: Различных разностей может быть 14 - от 1 до 14 - это те 14 клеток, в которые мы будем сажать кроликов. Кто же будет нашими кроликами? Ими, конечно, должны быть разности между парами данных нам натуральных чисел. Однако имеется 28 пар и их можно рассадить по 14 клеткам так, что в каждой клетке будет сидеть ровно два "кролика" (и значит, в каждой меньше трех). Здесь надо использовать дополнительное соображение: в клетке с номером 14 может сидеть не более одного кролика, ведь число 14 можно записать как разность двух натуральных чисел, не превосходящих 15, лишь одним способом: 14 = 15 - 1. Значит, в оставшихся 13 клетках сидят не менее 27 кроликов, и применение обобщенного принципа Дирихле дает нам желаемый результат.\\

\begin{flushleft}
  \textbf{\Large Обратный ход}
\end{flushleft}

\problem{1}Два пирата играли на золотые монеты. Сначала первый проиграл половину своих монет (отдал второму), потом второй проиграл половину своих, потом снова первый проиграл половину своих. В результате у первого оказалось 15 монет, а у второго — 33. Сколько монет было у первого пирата до начала игры?\\ 

\textbf{Ответ: 24.}\par
\textbf{Решение:} Попробуем проследить "с конца"\ ,  сколько у какого пирата было монет после каждой игры. В последней игре первый проиграл второму половину своих монет, 
\newpage
\noindent после чего у него осталось 15 монет. Но это ровно столько, сколько он только что отдал второму! Значит, перед этим у первого было $15 \cdot 2 = 30$ монет, а у второго было $33 - 15 = 18$ монет. Аналогично, во второй игре второй пират проиграл ровно столько, сколько у него осталось после второй игры, то есть 18 монет. Получаем, что перед второй игрой у второго пирата было $18 \cdot 2 = 36$ монет, а у первого было $30 - 18 = 12$ монет. Проделав самостоятельно ещё один шаг рассуждения, вы найдёте, что вначале у пиратов было по $24$ монеты.\\

\problem{2}Герман дал Эсти задание: "Сосчитай кляксы в своей тетради, прибавь к их числу 7, раздели на 8, умножь на 6 и отними 9. Если сделаешь всё правильно, получишь простое число". Эсти всё перепутал. Кляксы он подсчита точно, но потом умножил количество на 7, вычел из результата 8, затем разделил на 6 и прибавил 9. Какой ответ получился у Эсти?\\

\textbf{Ответ:} $18\dfrac{1}{6}$\par
\textbf{Решение:} Рассмотрим вычисления по плану Германа. Соседние операции "раздели на 8" \  и "умножь на 6" заменим на одну операцию умножения на $\dfrac{3}{4}$. Если бы после умножения на $\dfrac{3}{4}$ получалось дробное число, то, вычтя из него 9, мы бы снова получили дробное число, а должны получить простое (то есть целое), значит, при умножении на $\dfrac{3}{4}$ мы получаем целое число, причём кратное 3. Тогда после вычитания 9 получится число, также кратное 3. Единственное простое число, кратное 3, это само число 3. Значит, по плану Германа в конце должно было получиться 3. Произведя операции в обратном порядке, найдём число клякс: $(3 + 9) : 6 \cdot 8 - 7 = 9.$ Следовательно, Эсти получил число $(9 \cdot 7 - 8) : 6 + 9$ = $18\dfrac{1}{6}$\\

\begin{flushleft}
  \textbf{\Large Раскраска}
\end{flushleft}


\problem{1} Может ли конь, находясь на шахматной доске($8 \times 8$) в клетке h8, за 2024 хода попасть в клетку a8?\\

\textbf{Ответ: Нет, не может.}\par
\textbf{Решение:} Применим шахматную раскраску. Заметим, что после каждого своего
\newpage
\noindent хода конь меняет цвет, на котором стоит. Следовательно, после каждого второго хода(четного хода) конь будет на клетке того же цвета, на котором стоял. За 2024 хода он должен стоять на клетке того же цвета, с которой начинал, однако цвета клеток h8 и a8 разные. Значит, такого быть не может.\\

\problem{2}Дана доска размером $10\times10$. Можно ли разрезать доску на прямоугольники $1\times4$?\\

\textbf{Ответ: Нет, не можем.}\par
\textbf{Решение:} Применим диагональную раскраску в 4 цвета. Заметим, что где бы мы не вырезали прямоугольник 1x4, он всегда занимает клетки четырех разных цветов по одному. Так как всего клеток на доске $10\cdot10 = 100$, а прямоугольник занимает 4 клетки, значит, нам понадобится 100/4 = 25 прямоугольников, чтобы разрезать доску. То есть мы должны вырезать по 25 клеток каждого цвета. Однако на доске клеток первого цвета - 25, второго цвета - 26, третьего цвета - 25 и четвертого цвета - 24. Следовательно, у нас не получится разрезать доску на такие прямоугольники.\\
